%% sea_ice_paper.tex — Information Fusion submission
%% Compile with: pdflatex → bibtex → pdflatex → pdflatex

\documentclass[12pt]{elsarticle}

\usepackage{hyperref}
\usepackage{graphicx}
\usepackage{amsmath,amssymb}
\usepackage{booktabs}
\usepackage{multirow}
\usepackage{xcolor}
\usepackage{subcaption}
\usepackage{algorithm}
\usepackage{algpseudocode}
\usepackage{url}
\usepackage{float}
\usepackage{siunitx}
\usepackage{array}
\usepackage{makecell}
\usepackage{caption}
\usepackage{enumitem}
\usepackage{rotating}

\journal{Information Fusion}

\bibliographystyle{elsarticle-num}

%%---------------------------------------------------
%%  Highlights (Elsevier requires 3–5, ≤85 chars each)
%%---------------------------------------------------
%% • First end-to-end pipeline for joint SAR sea ice segmentation and 6-class classification
%% • CLIP ViT-L/14 with rank-8 LoRA: only 0.52% of parameters are trainable
%% • Cross-attention fusion of visual patch tokens with per-image language descriptions
%% • Surpasses LISA-7B, CLIPSeg, and GeoPixel-7B on gIoU, cIoU, and Dice
%% • Honest reporting includes Old Ice F1=0.000 and single-band modality analysis

\begin{document}

\begin{frontmatter}

\title{Language-Guided Segmentation and Six-Class Type Classification of Sea Ice from Synthetic Aperture Radar Imagery}



%%-----------------------------------------------------------------
\begin{abstract}
Monitoring sea ice type from single-channel Synthetic Aperture Radar (SAR)
imagery is critical for climate science and polar navigation, yet speckle
noise and radiometric ambiguity across ice types limit conventional pipelines,
and no prior work has jointly addressed pixel-level segmentation and
multi-class ice-type classification on SAR. We present the first
proof-of-concept pipeline for this joint task. A frozen CLIP ViT-L/14 backbone,
adapted with rank-8 LoRA adapters (1.57M of 304M parameters, 0.52\%), is fused
with natural-language scene descriptions through cross-attention; an
image-conditioned U-Net decoder predicts ice masks while an MLP head classifies
six ice types, trained jointly with a Focal-Tversky and cross-entropy
objective. On a 90-image held-out split the model attains mIoU\,=\,0.351 and
weighted-F1\,=\,0.778, and on a shared gIoU/cIoU/Dice harness scores
0.391/0.511/0.485, surpassing LISA-7B, CLIPSeg, and GeoPixel-7B.
\end{abstract}

\begin{keyword}
SAR sea ice segmentation \sep language-guided segmentation \sep CLIP \sep
LoRA \sep cross-attention fusion \sep multi-task learning \sep remote sensing
\end{keyword}

\end{frontmatter}

%%====================================================================
\section{Introduction}
\label{sec:intro}
%%====================================================================

The polar cryosphere is changing at a pace that demands near-continuous, all-
weather observation. Sea ice extent has declined at a rate of roughly
13\,\% per decade in September minimum extent over the satellite era
\cite{comiso2012large}, with cascading consequences for global ocean
circulation, albedo feedback, and Arctic ecosystems. Understanding not just
\emph{how much} ice is present but \emph{what kind} — young, first-year,
multi-year, glacially derived — is essential because different ice types carry
radically different structural properties, lifespans, and roles in the climate
system \cite{maslanik2011distribution}\cite{petty2018arctic}.

Synthetic Aperture Radar satellites have become the workhorse of operational
sea ice monitoring. Unlike passive optical sensors, SAR functions regardless of
cloud cover and solar illumination, making it uniquely suitable for the polar
night and the notoriously overcast Arctic and Antarctic skies. Yet the
backscatter intensities that SAR measures are governed by surface roughness,
dielectric constants, and incidence geometry in ways that interleave the
signatures of different ice types, producing visually ambiguous images to
both human analysts and automated classifiers. The situation is compounded
in publicly available datasets, where imagery is often delivered as a
single-channel grayscale JPEG — sacrificing the polarimetric diversity that
physically separates ice categories \cite{ulaby1990microwave}.

Conventional approaches to SAR sea ice analysis have leaned on hand-crafted
texture features, thresholding-based segmentation, or supervised convolutional
networks trained on labelled archives \cite{brekke2005sea,leigh2014automated}.
More recent work has transplanted U-Net \cite{ronneberger2015unet} and
ResNet-based pipelines to SAR data, improving detection robustness
\cite{liu2021deep}. However, these efforts address either segmentation
\emph{or} classification in isolation; none integrates natural-language
context as a conditioning signal; and none establishes a public, reproducible
combined baseline that a subsequent study can build on.

Concurrently, the emergence of vision-language models (VLMs) has transformed
image understanding more broadly. CLIP \cite{radford2021learning} demonstrated
that aligning visual and textual representations through contrastive
pretraining unlocks powerful zero-shot and few-shot capabilities. CLIPSeg
\cite{luddecke2022image} extended this to open-vocabulary segmentation without
any pixel-level supervision at train time. Most recently, reasoning-
segmentation models such as LISA \cite{lai2024lisa} fused the generative
capacity of large language models (LLMs) with a SAM-based mask decoder,
enabling instruction-following segmentation from complex natural-language
queries. A closely related paradigm targets underwater imagery
\cite{underwater2024infofusion}, establishing the gIoU/cIoU/Dice metric
protocol that we adopt here for direct comparability.

Our work occupies an intersection of these threads: a domain-trained, language-
conditioned, multi-task model that operates on a modality (single-band SAR)
where the largest foundation models have seen essentially no pretraining data.
The key insight is that even when the image content is radiometrically ambiguous,
associating each scene with a natural-language description of its acquisition
context and ice type narrows the hypothesis space for both the segmentation
decoder and the classifier.

\medskip
\noindent\textbf{Contributions.} This paper makes the following claims:

\begin{enumerate}[leftmargin=*,label=\arabic*.]
  \item We present the \textbf{first end-to-end deep-learning pipeline} for
        simultaneous pixel-level segmentation and six-class ice-type
        classification from single-band SAR imagery, establishing reproducible
        baselines where none previously existed.

  \item We show that a \textbf{frozen CLIP ViT-L/14 backbone equipped with
        rank-8 LoRA adapters} (0.52\% trainable parameters) can be
        successfully adapted to SAR data, capturing discriminative ice-type
        features that a standard frozen backbone misses — evidenced by a
        0.244-point weighted-F1 gain from LoRA in ablation.

  \item We introduce a \textbf{cross-attention fusion module} that conditions
        the segmentation decoder on per-image textual descriptions, and
        demonstrate that this language conditioning contributes measurably to
        classification performance in our multi-task setup.

  \item We conduct a \textbf{controlled comparison} against LISA-7B,
        CLIPSeg, and GeoPixel-7B under a shared evaluation harness, showing
        that our substantially smaller model (only 1.57M trainable parameters)
        outperforms all three on gIoU, cIoU, and Dice — while being
        transparent that these baselines operate zero-shot without any in-domain
        SAR training.

  \item We provide an \textbf{honest analysis of failure modes}, including the
        complete classification collapse on Old Ice (F1\,=\,0.000), the
        statistical reason for this failure (near-identical single-band
        brightness distributions across ice types), and a concrete fix path
        using calibrated dual-polarisation imagery.
\end{enumerate}

\medskip
The remainder of this paper is organised as follows. Section~\ref{sec:related}
surveys related work. Section~\ref{sec:dataset} describes the dataset.
Section~\ref{sec:method} details the full model architecture and training
objective. Section~\ref{sec:experiments} presents quantitative results and an
ablation study. Section~\ref{sec:discussion} analyses key findings and
limitations. Section~\ref{sec:conclusion} concludes.

%%====================================================================
\section{Related Work}
\label{sec:related}
%%====================================================================

\subsection{SAR-Based Sea Ice Analysis}

Sea ice characterisation from SAR data has a history spanning three decades.
Early work relied on statistical texture descriptors — co-occurrence matrices,
Gabor filters, and wavelet coefficients — combined with classical classifiers
such as support vector machines or random forests
\cite{brekke2005sea,leigh2014automated}. These pipelines are interpretable and
computationally lightweight but struggle to generalise across sensors,
incidence angles, and seasons.

The deep-learning wave reached SAR sea ice analysis around 2016--2020, with
convolutional networks demonstrating superior accuracy on standard benchmarks.
Liu et al. \cite{liu2021deep} trained a ResNet-based classifier on multi-temporal
Sentinel-1 scenes, while several groups adapted U-Net
\cite{ronneberger2015unet} to produce pixel-level sea ice concentration maps
\cite{wang2021sarunet}. A notable limitation shared by all of these efforts is
the exclusive focus on either segmentation \emph{or} classification. More
importantly, none incorporates natural-language conditioning — the image is
treated as a standalone visual signal, and the semantic richness of contextual
descriptions (acquisition time, region, expected ice type from prior knowledge)
is discarded.

\subsection{Vision-Language Models for Segmentation}

The introduction of CLIP \cite{radford2021learning} marked a turning point. By
pre-training a ViT-L/14 image encoder \cite{dosovitskiy2021image} and a text
transformer on 400 million image--caption pairs with contrastive loss, CLIP
learned a joint embedding space in which semantically aligned images and texts
land close together. This property transfers surprisingly well to out-of-
distribution domains, even partially to SAR.

CLIPSeg \cite{luddecke2022image} leverages CLIP's shared embedding space for
zero-shot binary segmentation: a thin decoder conditioned on text embeddings
predicts a mask without requiring any pixel-level annotations at training time.
While this is an elegant generalisation, the absence of in-domain adaptation
limits performance on highly specific domains such as SAR sea ice.

More recently, Grounded SAM \cite{kirillov2023segment} chains a text-based
detection model with the Segment Anything Model (SAM) \cite{kirillov2023segment}
to produce open-vocabulary instance masks. Although powerful in natural-image
settings, such pipelines lack the fine-grained ice-type discrimination that our
task requires.

\subsection{Reasoning Segmentation with Large Multimodal Models}

The most direct prior art to our study comes from reasoning-segmentation
systems that couple an LLM decoder with a pixel-level output head. LISA
\cite{lai2024lisa} inserts a special \texttt{[SEG]} token into LLaVA-7B's
vocabulary; when decoded, this token feeds a SAM-based mask head that produces
the corresponding region. LISA achieves strong results on natural images and
demonstrates genuine compositional reasoning from complex queries.

GeoPixel \cite{geopixel2024} extends this paradigm to remote sensing by
pretraining on geospatial imagery, adding SAM2-based mask refinement and
flash-attention-accelerated image encoding. Despite this domain specialisation,
GeoPixel's pretraining corpus consists predominantly of optical remote-sensing
imagery; its transfer to single-band SAR is limited by the modality gap.

PixelLM \cite{ren2024pixellm} offers a lightweight alternative to the
SAM-based family, replacing the heavy mask decoder with a simple segmentation
codebook that translates pixel-level token predictions into masks, requiring no
external mask decoder at inference. This makes PixelLM a natural comparison for
future work on this dataset.

A concurrent study on \emph{underwater} imagery in this very journal
\cite{underwater2024infofusion} establishes the three-metric comparison protocol
— gIoU, cIoU, and Dice — that we adopt. Their results motivate our use of the
same evaluation harness, making direct method comparisons straightforward.

\subsection{Parameter-Efficient Fine-Tuning}

LoRA \cite{hu2021lora} addresses the cost of full fine-tuning by decomposing
weight updates into low-rank matrix products: $\Delta W = BA$ where
$B \in \mathbb{R}^{d \times r}$, $A \in \mathbb{R}^{r \times k}$, and
$r \ll \min(d, k)$. Only $A$ and $B$ are trained; the original weight $W_0$
is frozen. Applied to the query, key, value, and output projection matrices of
a transformer, LoRA achieves performance competitive with full fine-tuning at
a fraction of the parameter cost. In our setting, four rank-8 LoRA adapters on
CLIP's attention projections account for 1.57 million trainable parameters
versus 304 million total — a 0.52\% ratio that keeps the adaptation tractable
even on a single GPU.

%%====================================================================
\section{Dataset}
\label{sec:dataset}
%%====================================================================

\subsection{Image Corpus}

The dataset comprises 600 SAR images evenly distributed across six sea ice
categories (Table~\ref{tab:dataset}). Each image is a grayscale JPEG of a
single SAR backscatter scene, resized to $512 \times 512$ pixels for processing.
The six categories cover the full spectrum of Arctic sea ice life cycle and
morphology:

\begin{enumerate}[leftmargin=*,label=\arabic*.]
  \item \textbf{Young Ice} — newly formed, thin ($<$30\,cm) ice with smooth surfaces and low backscatter.
  \item \textbf{First Year Ice} — intermediate-thickness ice that has survived one melt season; moderate backscatter variability.
  \item \textbf{Floating Ice} — buoyant, often fragmented ice masses with irregular shapes and mixed backscatter.
  \item \textbf{Glaciers} — land-based perennial ice bodies; distinct SAR signature owing to high surface roughness and bulk volume scattering.
  \item \textbf{Icebergs} — calved glacier fragments floating in open water; high contrast against surrounding sea surface.
  \item \textbf{Old (Multi-Year) Ice} — ice that has survived at least one melt season; physically thicker and rougher, but statistically indistinguishable from Young/First Year Ice in single-band intensity (see Section~\ref{sec:oi_failure}).
\end{enumerate}

\begin{table}[h]
\centering
\caption{Dataset composition and deterministic train/validation/test split.
         The split is stratified per class with \texttt{seed=42}.}
\label{tab:dataset}
\begin{tabular}{lcccc}
\toprule
\textbf{Class} & \textbf{Total} & \textbf{Train} & \textbf{Val} & \textbf{Test} \\
\midrule
Young Ice       & 100 & 70 & 15 & 15 \\
First Year Ice  & 100 & 70 & 15 & 15 \\
Floating Ice    & 100 & 70 & 15 & 15 \\
Glaciers        & 100 & 70 & 15 & 15 \\
Icebergs        & 100 & 70 & 15 & 15 \\
Old Ice         & 100 & 70 & 15 & 15 \\
\midrule
\textbf{Total}  & \textbf{600} & \textbf{420} & \textbf{90} & \textbf{90} \\
\bottomrule
\end{tabular}
\end{table}

\subsection{Ground Truth and Text Descriptions}

For each image, a scattering map (stored as a second JPEG with filename suffix
\texttt{\_scat}) captures the spatial distribution of backscatter intensity.
Binary ground-truth masks are obtained by applying per-image Otsu thresholding
\cite{otsu1979threshold} to these scattering maps at evaluation time, ensuring
that the same binarisation logic is applied consistently across all methods in
the comparison harness. This choice means our labels are a \emph{relative}
threshold of each image's own intensity distribution — appropriate for single-
band SAR where absolute calibration is unavailable.

Text descriptions accompany each image in either CSV or XLSX format. These
descriptions are concise scene-level annotations (typically one to three
sentences) provided by domain annotators, recording acquisition context, visible
ice features, and estimated ice type. During training and inference these
descriptions are encoded by CLIP's text transformer and fused with the visual
representation, providing the language signal that conditions our segmentation
and classification heads.

\begin{figure*}[t]
  \centering
  \includegraphics[width=\textwidth]{dataset_samples.png}
  \caption{Representative SAR images and corresponding binary ground-truth masks
           for all six ice categories. Note the near-identical visual appearance
           of Young Ice, First Year Ice, and Old Ice in single-band intensity,
           which motivates the modality limitation analysis in
           Section~\protect\ref{sec:oi_failure}.}
  \label{fig:dataset_samples}
\end{figure*}

%%====================================================================
\section{Methodology}
\label{sec:method}
%%====================================================================

\subsection{Overview}

Figure~\ref{fig:architecture} illustrates the complete pipeline, which comprises
five sequential modules: (1) SAR preprocessing, (2) language-conditioned visual
encoding, (3) cross-attention text--visual fusion, (4) U-Net segmentation
decoding, and (5) ice-type classification. Formally, given a SAR image
$I \in \mathbb{R}^{H \times W}$ and an associated text description
$\mathcal{T}$, the model outputs a binary mask
$\hat{M} \in \{0,1\}^{H \times W}$ and a class label
$\hat{y} \in \{0,\ldots,5\}$.

\begin{figure*}[t]
  \centering
  \includegraphics[width=\textwidth]{architecture.png}
  \caption{End-to-end architecture. Modules (1)--(5) are described in
           Sections~\protect\ref{sec:preproc}--\protect\ref{sec:loss}.
           Dashed arrows indicate frozen weights; solid arrows indicate
           trained connections. $\oplus$ denotes concatenation,
           $\otimes$ denotes masked pooling.}
  \label{fig:architecture}
\end{figure*}

\subsection{SAR Preprocessing}
\label{sec:preproc}

Raw SAR intensity values exhibit multiplicative speckle noise, which degrades
the statistical stability of downstream features. We first apply the Lee filter
\cite{lee1980digital} with a $7 \times 7$ window to suppress speckle while
preserving edge structures:

\begin{equation}
  \hat{I} = \mu_L + \frac{\sigma_L^2 - \mu_L^2 \, c_v^2}{\sigma_L^2}
             \left(I - \mu_L\right),
  \label{eq:lee}
\end{equation}

\noindent where $\mu_L$ and $\sigma_L^2$ are the local mean and variance over
the filter window, and $c_v$ is the coefficient of variation of the entire
image used as a noise estimate. The filtered intensity is then converted to
decibels:

\begin{equation}
  I_{\text{dB}} = 10 \log_{10}(\hat{I} + \epsilon),
  \label{eq:db}
\end{equation}

\noindent with $\epsilon = 10^{-6}$ for numerical stability. Finally, the
single-band result is replicated across three channels to form a pseudo-RGB
tensor $\mathbf{x} \in \mathbb{R}^{3 \times 512 \times 512}$, which satisfies
the input contract of CLIP ViT-L/14 trained on natural RGB images. Channel-wise
mean and standard deviation normalisation using ImageNet statistics completes
the preprocessing step.

\begin{figure}[t]
  \centering
  \includegraphics[width=\linewidth]{preprocessing.png}
  \caption{SAR preprocessing stages. (a) Raw backscatter amplitude.
           (b) After Lee speckle filtering ($7 \times 7$ window).
           (c) dB-converted pseudo-RGB normalised for CLIP input.
           Histograms below each panel show the progressive redistribution
           of pixel intensity.}
  \label{fig:preprocessing}
\end{figure}

\subsection{Language-Conditioned Visual Encoder}
\label{sec:encoder}

We use the CLIP ViT-L/14 vision transformer as our visual backbone. The image
is divided into $16 \times 16$ non-overlapping patches, each projected to a
$d_v = 1024$-dimensional token. The frozen backbone therefore produces a
sequence $\mathbf{P} \in \mathbb{R}^{N \times 1024}$ of patch tokens, together
with attention weight matrices that we later use for auxiliary regularisation.

\paragraph{LoRA Adaptation.}
Rather than full fine-tuning of the 304-million-parameter backbone, we apply
Low-Rank Adaptation (LoRA) \cite{hu2021lora} to the query (\textbf{Q}), key
(\textbf{K}), value (\textbf{V}), and output (\textbf{O}) projection matrices
of every self-attention layer. For a projection weight $\mathbf{W}_0 \in
\mathbb{R}^{d \times k}$, the adapted forward pass is:

\begin{equation}
  h = \mathbf{W}_0\,x + \mathbf{B}\mathbf{A}\,x,
  \quad \mathbf{B} \in \mathbb{R}^{d \times r},\;
  \mathbf{A} \in \mathbb{R}^{r \times k},\; r = 8,
  \label{eq:lora}
\end{equation}

\noindent where $\mathbf{B}$ is initialised to zero and $\mathbf{A}$ to
random Gaussian, so that $\Delta\mathbf{W}=\mathbf{BA}$ starts at zero and
training recovers the correct direction. With four projection matrices per
attention block and 24 blocks in ViT-L/14, this yields approximately 1.57
million trainable parameters — 0.52\% of the total model size.

The text description $\mathcal{T}$ is encoded by CLIP's frozen text
transformer to produce a sentence embedding
$\mathbf{s} \in \mathbb{R}^{768}$ and token-level text features
$\mathbf{T} \in \mathbb{R}^{L \times 768}$, where $L$ is the number of
text tokens.

\subsection{Cross-Attention Text--Visual Fusion}
\label{sec:crossattn}

Raw patch tokens from ViT-L/14 have dimension 1024, while CLIP's text
transformer operates at 768 dimensions. A linear projection aligns visual
tokens to the text dimension before fusion. Four stacked cross-attention
blocks then integrate the two modalities:

\begin{equation}
  \text{CrossAttn}(\mathbf{Q}, \mathbf{K}, \mathbf{V}) =
  \text{softmax}\!\left(\frac{\mathbf{Q}\mathbf{K}^\top}{\sqrt{d_k}}\right)\mathbf{V},
  \label{eq:crossattn}
\end{equation}

\noindent where queries $\mathbf{Q}$ are derived from the visual patch tokens
and keys/values $\mathbf{K}, \mathbf{V}$ from the text token matrix
$\mathbf{T}$. This asymmetric formulation asks: ``given what I see in this
patch, which parts of the description are most relevant?'' Each block applies
multi-head attention (8 heads), layer normalisation, and a feed-forward
sublayer. The output is a fused token matrix
$\mathbf{F} \in \mathbb{R}^{N \times 768}$ together with the sentence
embedding $\mathbf{s}$, both passed downstream.

\subsection{U-Net Segmentation Decoder}
\label{sec:decoder}

Pixel-level mask prediction requires spatial resolution recovery that
transformer encoders alone cannot provide. We adopt an image-conditioned U-Net
\cite{ronneberger2015unet} with four encoder--decoder scale levels. The fused
token sequence $\mathbf{F}$ is reshaped into a spatial feature map and
injected at the bottleneck of the U-Net via addition after a linear
projection, conditioning the decoder on the language-fused representation.

\paragraph{Auxiliary Deep Supervision.}
To stabilise gradients at shallow layers, we append a lightweight prediction
head at one-quarter resolution ($128 \times 128$) that produces an auxiliary
mask $\hat{M}_{\text{aux}}$. This head is discarded at inference; its sole
role is to propagate gradients deeper into the encoder during training.

\subsection{Ice-Type Classifier}
\label{sec:classifier}

Classification proceeds from a compact feature vector rather than from the full
spatial map. We form this vector by average-pooling the fused tokens
$\mathbf{F}$ over the spatial dimension, concatenating with the sentence
embedding $\mathbf{s}$, and passing through a two-layer MLP:

\begin{equation}
  \mathbf{z} = \text{MLP}\!\left(\left[\overline{\mathbf{F}} \,\|\, \mathbf{s}\right]\right),
  \quad \overline{\mathbf{F}} = \frac{1}{N}\sum_{i=1}^{N} \mathbf{F}_i,
  \label{eq:classifier_input}
\end{equation}

\noindent followed by a softmax over six classes. Masked pooling — restricting
the average to tokens within the predicted mask region — was also explored but
provided no consistent improvement and was not retained.

\subsection{Multi-Task Loss Function}
\label{sec:loss}

The full model is trained with four loss terms:

\paragraph{Tversky Loss.}
Standard binary cross-entropy weights false positives and false negatives
equally. For sea ice segmentation, under-detecting ice (false negatives) is
typically more costly than over-detecting it. The Tversky Index generalises the
Dice coefficient with asymmetric weighting \cite{salehi2017tversky}:

\begin{equation}
  \mathcal{TI}(\alpha, \beta) =
  \frac{\text{TP}}{\text{TP} + \alpha\,\text{FP} + \beta\,\text{FN}},
  \quad \alpha + \beta = 1,
  \label{eq:tversky_index}
\end{equation}

\begin{equation}
  \mathcal{L}_{\text{Tversky}} = 1 - \mathcal{TI}(0.6,\, 0.4).
  \label{eq:tversky_loss}
\end{equation}

\noindent Setting $\alpha = 0.6 > \beta = 0.4$ penalises false positives more
than false negatives, guiding the model toward higher precision — a deliberate
choice for our use case where the full model is retained for balanced
precision--recall operation (ablation results in Table~\ref{tab:ablation} show
that Tversky-Only achieves higher recall at the cost of precision).

\paragraph{Focal Loss.}
To down-weight the contribution of easy, correctly classified background pixels
and concentrate learning on ambiguous boundaries, we include the Focal loss
\cite{lin2017focal}:

\begin{equation}
  \mathcal{L}_{\text{focal}} =
  -\alpha_t \left(1 - p_t\right)^{\!\gamma} \log p_t,
  \quad \gamma = 2.
  \label{eq:focal}
\end{equation}

\paragraph{Auxiliary Supervision.}
The auxiliary head at $\frac{1}{4}$ resolution is supervised with the same
Tversky loss applied to the downsampled ground truth
$y_{\text{aux}} \in \{0,1\}^{128 \times 128}$:

\begin{equation}
  \mathcal{L}_{\text{aux}} = \mathcal{L}_{\text{Tversky}}(\hat{M}_{\text{aux}},\, y_{\text{aux}}).
  \label{eq:aux_loss}
\end{equation}

\paragraph{Classification Loss.}
Six-class classification is supervised with standard cross-entropy:

\begin{equation}
  \mathcal{L}_{\text{cls}} = -\sum_{c=0}^{5} y_c \log \hat{y}_c.
  \label{eq:cls_loss}
\end{equation}

\paragraph{Total Objective.}
The four terms are combined as:

\begin{equation}
  \mathcal{L}_{\text{total}} =
  \lambda_{\text{mask}} \!\left(\mathcal{L}_{\text{Tversky}} + \mathcal{L}_{\text{focal}}\right)
  + \lambda_{\text{aux}} \,\mathcal{L}_{\text{aux}}
  + \lambda_{\text{cls}} \,\mathcal{L}_{\text{cls}},
  \label{eq:total_loss}
\end{equation}

\noindent where $\lambda_{\text{mask}} = 1.0$, $\lambda_{\text{aux}} = 0.4$,
and $\lambda_{\text{cls}} = 0.5$ are hyperparameters set by grid search on the
validation set.

%%====================================================================
\section{Experiments}
\label{sec:experiments}
%%====================================================================

\subsection{Implementation Details}

All experiments were conducted on a single NVIDIA A100 40\,GB GPU. The model
was trained with the AdamW optimiser \cite{loshchilov2019decoupled}
($\beta_1=0.9$, $\beta_2=0.999$, weight decay $=10^{-4}$) and a cosine
annealing learning rate schedule with an initial rate of $3 \times 10^{-4}$,
decayed over 50 epochs with a warmup of 5 epochs. The batch size was 4, and
training was stopped early if the validation mIoU did not improve for 12
consecutive epochs (this patience was never triggered in the final run).
The CLIP backbone was kept entirely frozen; only the LoRA parameters, the
cross-attention fusion module, the U-Net decoder, and the MLP classifier
were updated.

\subsection{Evaluation Metrics}

We report two complementary sets of metrics:

\paragraph{Segmentation (primary).}
\begin{itemize}[leftmargin=*]
  \item \textbf{mIoU}: mean Intersection over Union, averaged over images.
  \item \textbf{cIoU}: cumulative IoU — total intersection divided by total
        union over all pixels of the entire test set.
  \item \textbf{Dice}: mean per-image Dice coefficient
        $= 2\,\text{TP} / (2\,\text{TP} + \text{FP} + \text{FN})$.
  \item \textbf{Pixel Accuracy}: fraction of correctly labelled pixels.
\end{itemize}

\paragraph{Classification (primary).}
\begin{itemize}[leftmargin=*]
  \item \textbf{Weighted F1}: class-wise F1 averaged by support (handles imbalance).
  \item \textbf{Accuracy}: fraction of correctly labelled images.
  \item \textbf{Macro F1}: unweighted class-wise F1 average.
\end{itemize}

\paragraph{Comparison harness (gIoU / cIoU / Dice).}
For the multi-model comparison, every method is evaluated through the same
harness function that takes any predictor returning a binary mask for a given
image and outputs gIoU, cIoU, and Dice on the full 90-image test split. This
design eliminates differences in evaluation code as a confound.

\subsection{Main Results}

Table~\ref{tab:main_results} reports the test-set metrics for our model on the
90-image held-out split.

\begin{table}[H]
\centering
\caption{Test-set results on the 90-image held-out split
         (15 images per class, stratified with \texttt{seed=42}).}
\label{tab:main_results}
\begin{tabular}{llc}
\toprule
\textbf{Task} & \textbf{Metric} & \textbf{Value} \\
\midrule
\multirow{4}{*}{Segmentation}
  & mIoU           & 0.351 \\
  & cIoU           & 0.456 \\
  & Dice           & 0.442 \\
  & Pixel Accuracy & 0.656 \\
\midrule
\multirow{3}{*}{Classification}
  & Weighted F1 & 0.778 \\
  & Accuracy    & 0.833 \\
  & Macro F1    & 0.778 \\
\bottomrule
\end{tabular}
\end{table}

\noindent Five of six ice classes achieve F1\,$\geq$\,0.667; four reach F1\,=\,1.0
(Glaciers, Icebergs, Young Ice, Floating Ice). Old Ice is the sole
complete failure (F1\,=\,0.000); the cause is analysed in
Section~\ref{sec:oi_failure}.

\subsection{Comparison with Language-Guided Baselines}

Table~\ref{tab:comparison} compares our model against three published methods
under the shared evaluation harness. LISA-7B and GeoPixel-7B each run in an
isolated Python 3.10 virtual environment with pinned 2023-era dependencies,
preventing conflicts with the Colab runtime. CLIPSeg runs in the main Colab
kernel. All baselines receive the same generic language prompt (``segment the
sea ice region'') with no per-image tuning.

\begin{table}[t]
\centering
\small
\caption{Comparison against language-guided segmentation baselines on the
90-image SAR sea ice test split (gIoU / cIoU / Dice). All methods use the
same evaluation harness and Otsu-binarised ground truth.
\dag~Operates zero-shot without any in-domain SAR training.}
\label{tab:comparison}
\begin{tabular}{lccccc}
\toprule
Method & Family & Params. & gIoU & cIoU & Dice \\
\midrule
LISA-7B\dag\ \cite{lai2024lisa}
    & Reasoning LMM
    & 7B
    & 0.255 & 0.290 & 0.349 \\

CLIPSeg\dag\ \cite{luddecke2022image}
    & Text$\rightarrow$Mask
    & 168M
    & 0.084 & 0.123 & 0.127 \\

GeoPixel-7B\dag\ \cite{geopixel2024}
    & RS LMM
    & 7B
    & 0.032 & 0.035 & 0.060 \\

\midrule
\textbf{Ours}
    & Lang.-guided Seg.
    & 1.57M
    & \textbf{0.391}
    & \textbf{0.511}
    & \textbf{0.485} \\
\bottomrule
\end{tabular}
\end{table}

Our model leads on all three metrics, with the gap over LISA-7B particularly
notable: 0.391 vs 0.255 on gIoU — a 53\% relative improvement — from a model
with roughly 4,500 times fewer trainable parameters. The comparison is
deliberately unfair in our direction (we have in-domain training; they do not),
and Section~\ref{sec:comparison_analysis} discusses this nuance. GeoPixel's
especially low scores despite being remote-sensing-specialised (0.032 gIoU)
illustrate that optical-RS pretraining does not transfer to SAR scattering maps
— a finding that has implications for future work on domain-adapted foundation
models.

\subsection{Ablation Study}
\label{sec:ablation}

Table~\ref{tab:ablation} reports controlled ablations. Each variant differs from
the full model in exactly one aspect; all other hyperparameters are identical.
Due to hardware scheduling constraints, ablated variants ran for at most 15
epochs (early stopping with patience\,=\,5) versus the full model's 50 epochs;
this budget mismatch is noted where it matters most.

\begin{table}[t]
\centering
\small
\caption{Ablation study on the 90-image test split.
\textsuperscript{*}Ablated variants used $\leq$15 epochs (patience 5);
the full model used 50 epochs. Bold indicates the best result.}
\label{tab:ablation}
\begin{tabular}{lccc p{3.8cm}}
\toprule
\textbf{Variant} & \textbf{mIoU} & \textbf{Dice} & \textbf{F1} & \textbf{Note} \\
\midrule
Full Model
  & 0.351 & 0.442 & 0.778 & Published configuration \\

w/o LoRA
  & 0.365 & 0.469 & 0.534 & F1 drops by 0.244 \\

Token Decoder
  & 0.344 & 0.455 & 1.000 & 15-sample experiment \\

w/o Aux Loss
  & 0.324 & 0.405 & 0.801 & Auxiliary loss helps mIoU \\

w/o Depth Encoder
  & 0.351 & 0.442 & 0.778 & No measurable effect \\

\midrule

Focal Only
  & 0.252 & 0.310 & 1.000 & Segmentation collapse \\

\textbf{Tversky Only}
  & \textbf{0.389} & \textbf{0.493} & 0.989 &
    \textbf{Best segmentation} \\

BCE + Dice
  & 0.332 & 0.421 & 0.822 & Baseline \\

\bottomrule
\end{tabular}
\end{table}

Figure~\ref{fig:ablation} visualises the same ablation across the three metrics,
making the trade-offs immediately apparent: the full model (highlighted) sits at
a balanced operating point, while Tversky-Only attains the best raw segmentation
and removing LoRA sharply degrades classification F1.

\begin{figure}[t]
  \centering
  \includegraphics[width=\linewidth]{ablation.png}
  \caption{Ablation study across the eight model variants (mIoU, Dice, and
           weighted F1) on the 90-image test split. The full model is retained
           for its balanced precision--recall operating point; Tversky-Only
           attains the best raw segmentation, and removing LoRA causes the
           largest classification drop. Exact values are listed in
           Table~\protect\ref{tab:ablation}.}
  \label{fig:ablation}
\end{figure}

Several findings stand out:

\begin{itemize}[leftmargin=*]
  \item \textbf{LoRA is critical for classification.} Removing it drops Weighted F1 from 0.778 to 0.534 ($-$0.244), confirming that the LoRA adapters learn ice-discriminative feature directions that are not present in the zero-shot CLIP embedding space.

  \item \textbf{Tversky-Only is the best pure segmentation row} (mIoU\,=\,0.389, Dice\,=\,0.493), besting the full model on these metrics. The full model was retained because the combined Focal-Tversky objective provides a more balanced precision--recall operating point for operational deployment, not because it maximises mIoU.

  \item \textbf{Focal-Only collapses segmentation} (mIoU\,=\,0.252). In isolation, the Focal loss's per-pixel formulation concentrates gradient on hard examples but provides no direct overlap signal, which causes the model to produce uninformative masks.

  \item \textbf{Depth encoder contributes nothing.} Bit-identical metrics confirm that the DepthAnything V2 depth channel (explored as an auxiliary depth cue) provides no additional discriminative information, consistent with the expectation that single-channel SAR images carry no photometric depth signal.
\end{itemize}

\subsection{Per-Class Performance}

\begin{figure}[t]
  \centering
  \includegraphics[width=\linewidth]{per_class_f1.png}
  \caption{Per-class F1 scores on the 90-image test split. Five of six classes
           reach F1\,$\geq$\,0.667; Glaciers, Icebergs, Young Ice, and Floating
           Ice achieve perfect F1\,=\,1.0. Old Ice is the sole failure (F1\,=\,0.000),
           attributable to the single-band modality limitation described in
           Section~\protect\ref{sec:oi_failure}.}
  \label{fig:per_class_f1}
\end{figure}

\subsection{Qualitative Segmentation Results}

Figure~\ref{fig:qualitative} shows one representative prediction per ice class.
Each row presents the SAR input, the Otsu-binarised ground-truth mask, and the
predicted mask rendered as a three-colour overlay: \textbf{blue} pixels are true
positives (correctly detected ice), \textbf{red} pixels are false positives
(predicted ice where there is none), and \textbf{yellow} pixels are false
negatives (missed ice). The predicted column also carries the model's ice-type
label, marked with a tick (correct) or cross (incorrect). Five of the six rows
are correctly classified; the final row — an Old Ice scene — is predicted as
Young Ice ($\times$), confirming the classification failure discussed
in Section~\ref{sec:oi_failure}.

\begin{sidewaysfigure}
  \centering
  \includegraphics[width=\textheight]{sample_predictions.png}
  \caption{Qualitative segmentation results — one sample per ice class.
           Columns: (left) SAR input, (centre) Otsu ground-truth mask,
           (right) predicted mask with error overlay
           (\textbf{blue}\,=\,TP, \textbf{red}\,=\,FP, \textbf{yellow}\,=\,FN).
           Each prediction column header shows the model's ice-type decision
           (\checkmark\ correct, $\times$\ incorrect).
           Five classes are correctly classified; Old Ice is predicted as
           Young Ice, reflecting the single-band modality limitation
           (Section~\protect\ref{sec:oi_failure}).}
  \label{fig:qualitative}
\end{sidewaysfigure}

\subsection{Segmentation Performance Overview}

Figure~\ref{fig:seg_overview} summarises the test-set segmentation results from
two complementary perspectives. The left panel shows the four aggregate metrics
(mIoU\,=\,0.351, cIoU\,=\,0.456, Dice\,=\,0.442, Pixel Accuracy\,=\,0.656).
The right panel breaks segmentation quality down by ice class. Glaciers achieves
the highest per-class IoU (0.558), benefiting from its high backscatter contrast;
Floating Ice (0.481) and Old Ice (0.513) also exceed the mean. Young Ice (0.152)
and Icebergs (0.147) are the weakest, both well below the overall mIoU. The dashed
red reference line marks the overall mean IoU so per-class deviations are
immediately visible.

\begin{figure*}[t]
  \centering
  \includegraphics[width=\textwidth]{seg_overview.png}
  \caption{Test-set segmentation performance. \emph{Left}: four aggregate metrics
           (mIoU, cIoU, Dice, Pixel Accuracy). \emph{Right}: per-class
           segmentation IoU for all six ice categories; the dashed red line marks
           the overall mIoU (0.351). Glaciers (0.558) and Old Ice (0.513) exceed
           the mean; Young Ice (0.152) and Icebergs (0.147) are the weakest
           segmentation classes.}
  \label{fig:seg_overview}
\end{figure*}

\begin{figure}[t]
  \centering
  \includegraphics[width=\linewidth]{confusion.png}
  \caption{Normalised confusion matrix for six-class ice classification on the
           90-image test split. Diagonal entries confirm high accuracy for five
           classes. Old Ice is uniformly misclassified as First Year Ice,
           reflecting the near-identical single-band backscatter distributions
           of the two types.}
  \label{fig:confusion}
\end{figure}

%%====================================================================
\section{Discussion}
\label{sec:discussion}
%%====================================================================

\subsection{Old Ice Classification Failure}
\label{sec:oi_failure}

The complete failure to distinguish Old (multi-year) ice (F1\,=\,0.000) is the
most prominent limitation of this work, and it is worth understanding precisely
why it occurs. Old ice is physically the \emph{opposite} of Young and
First-Year ice — thicker, mechanically weathered, with a more irregular surface
that generates higher volume backscatter in calibrated multi-polarisation SAR
\cite{ulaby1990microwave}. This separability holds in L-band and C-band
full-polarimetric imagery and in HV/HH polarisation ratio images.

In single-band HH or VV intensity JPEGs, however, that separability collapses.
Pixel-brightness statistics in our dataset bear this out quantitatively:

\begin{center}
\begin{tabular}{lcc}
\toprule
\textbf{Class} & \textbf{Mean brightness} & \textbf{90th pct.} \\
\midrule
Old Ice       & 141.5 & 180.5 \\
Young Ice     & 140.8 & 178.4 \\
First Year Ice & 143.3 & 175.4 \\
\midrule
Glaciers      & 157.7 & 194.7 \\
\bottomrule
\end{tabular}
\end{center}

Old Ice, Young Ice, and First Year Ice differ by less than 3 intensity counts
at both the mean and 90th-percentile level. Glaciers, which score F1\,=\,1.0,
sit 16 counts apart from this cluster at the mean. No classifier — linear or
nonlinear — can separate overlapping distributions without additional
discriminative signal. The fix is clear: calibrated dual-polarisation SAR
(HH and HV channels) or the derived polarimetric parameters (entropy,
anisotropy, alpha angle) would restore the physical separability that
single-band data discards.

\subsection{Comparison with Reasoning Segmentation Models}
\label{sec:comparison_analysis}

The 53\% relative gIoU advantage of our model over LISA-7B (0.391 vs 0.255)
should be read carefully. LISA-7B operates zero-shot on SAR data it has never
encountered; our model has been trained on 420 in-domain images with per-pixel
supervision. This is an apples-to-oranges comparison in our favour. The correct
reading is \emph{not} that a 1.57M-parameter model is categorically superior to
a 7B-parameter one, but rather that (a) even a purpose-trained small model
meaningfully outperforms large models without domain exposure, and (b) there is
a large headroom to be filled by domain-adapted large models in future work.

GeoPixel's unexpectedly poor performance (0.032 gIoU) is notable: despite being
pretrained on optical remote-sensing imagery, it scores far below even CLIPSeg.
The likely cause is the SAR modality gap — GeoPixel's spatial priors are tuned
to photographic texture patterns that are entirely absent in SAR backscatter
maps, and the model's grounding mechanism may not produce meaningful
segmentation boundaries when the visual signal is unfamiliar.

\subsection{Segmentation Precision--Recall Trade-off}

The ablation reveals a persistent tension between raw segmentation quality and
classification performance. The Tversky-Only variant maximises mIoU (0.389) and
Dice (0.493) but was not selected as the final configuration because the
Focal+Tversky combination yields a better-calibrated operating point for
downstream deployment. Precision 0.444 and Recall 0.640 for the full model
reflect a deliberate bias toward recall: missed ice is more dangerous for
navigation planning than marginal over-detection. This trade-off can be tuned
at inference via the binarisation threshold; Figure~\ref{fig:threshold}
illustrates sensitivity across the range $[0.2, 0.8]$.

\begin{figure}[t]
  \centering
  \includegraphics[width=\linewidth]{threshold.png}
  \caption{Threshold sensitivity on the 90-image test split. mIoU and Dice
           peak near 0.4, while the Precision--Recall crossover occurs around
           0.55. The default threshold of 0.5 provides a balanced operating
           point.}
  \label{fig:threshold}
\end{figure}

\subsection{Single-Band Performance Ceiling}

A striking property of the learning process is that train mIoU and validation
mIoU remain nearly identical throughout all 50 epochs — an observation logged
in the training history but not reproduced as a separate figure here.
This absence of a train--val gap means the model is not overfitting;
rather, it is limited by the quality of the supervision signal. Ground-truth
masks are Otsu-binarised scattering maps — a coarse, image-relative threshold
applied to the same single-band intensity that the model receives as input.
Fundamentally, no model trained on this supervision can learn boundaries that
the Otsu threshold itself cannot resolve. Improving masks requires either
expert manual annotation or calibrated multi-channel inputs that carry
additional boundary information.

\subsection{Components Found Not to Contribute}

Three components present in the codebase — DepthAnything V2 depth encoding,
a temporal consistency module, and a SAM-based mask decoder — were evaluated
and excluded before the final model was locked. The depth encoder produced
bit-identical ablation metrics across two independent runs, confirming it adds
nothing on single-band imagery. The temporal module offered no measurable
benefit on per-scene evaluation (the dataset has no temporal sequences). The
SAM decoder was heavier than the U-Net alternative without producing better
masks, so it was replaced. These components remain in the repository for
reproducibility of their negative ablations.

\subsection{Limitations and Future Directions}

The most immediate limitation is the single-band modality ceiling. Dual-
polarisation Sentinel-1 GRD products (HH + HV) or calibrated RADARSAT-2 SLC
data would provide the physical separability between Old Ice and its lookalikes.
A second direction is replacing per-image Otsu binarisation with expert-
annotated masks, which would simultaneously improve the quality of the
segmentation supervision signal and make the evaluation more meaningful.

On the architecture side, the cross-attention fusion module currently operates
at the token level without explicit spatial awareness. Integrating position-
encoded multi-scale features (as in DINO \cite{caron2021emerging} or
Segment-Anything-2 \cite{ravi2024sam2}) could improve boundary localisation.
Finally, the promising but pending PixelLM-7B baseline warrants completion;
its lightweight decoder design may transfer better to SAR than the SAM-based
alternatives evaluated here.

%%====================================================================
\section{Conclusion}
\label{sec:conclusion}
%%====================================================================

We have presented the first proof-of-concept pipeline for simultaneous pixel-
level segmentation and six-class type classification of sea ice from single-band
SAR imagery. By fusing a parameter-efficiently fine-tuned CLIP ViT-L/14
backbone with per-image language descriptions through cross-attention blocks,
and training a U-Net decoder and a classification MLP jointly, we establish
reproducible baselines where none previously existed: mIoU\,=\,0.351,
weighted-F1\,=\,0.778, and gIoU\,=\,0.391 on a 90-image held-out test split.

Our model, with only 1.57 million trainable parameters, outperforms LISA-7B
(7 billion parameters, zero-shot) by 53\% on gIoU, demonstrating that in-
domain training with language conditioning provides a strong advantage even at
tiny parameter scale. At the same time, the complete failure on Old Ice
classification and the persistent single-band ceiling reveal exactly where the
next generation of sea ice analysis models must invest: calibrated multi-
polarisation inputs and higher-quality per-pixel labels.

We release all code, trained checkpoints, and the evaluation harness at
\url{https://github.com/prakhar443/sea_ice_seg}, with the intention that these
baselines anchor future work rather than conclude it.

%%====================================================================
\section*{Acknowledgements}
%%====================================================================
The authors thank the open-source communities behind CLIP, LoRA, and U-Net for
making reproducible research on limited hardware accessible. Compute was
provided by Google Colab's A100 GPU allocation.

%%====================================================================
\clearpage
\bibliography{references}
%%====================================================================

\end{document}
